\documentclass[11pt,a4paper]{article}
\usepackage[T1]{fontenc}
\usepackage[utf8]{inputenc}
\usepackage[main=italian,provide=*]{babel}
\usepackage{amsmath,amssymb}
\usepackage{mathtools}
\usepackage{xcolor}
\usepackage[a4paper,margin=2.7cm]{geometry}
\usepackage{caption}
\usepackage{booktabs}
\usepackage{enumitem}
\usepackage{placeins}
\usepackage[colorlinks=true,linkcolor=black,citecolor=black,urlcolor=blue!55!black]{hyperref}
\newcommand{\ket}[1]{\lvert #1 \rangle}
\newcommand{\bra}[1]{\langle #1 \rvert}

\title{Guida d'uso --- pacchetto riproducibile\\Parte 1 (sistema chiuso), trimero ad anello}
\author{Tesi triennale, Samuele --- Relatore: Prof.\ Alessandro Chiesa}
\date{}

\begin{document}
\sloppy
\maketitle

\tableofcontents

\section{Scope di questo pacchetto: cosa copre, cosa no}

Copre la pipeline computazionale canonica di Parte 1 sull'anello:
diagonalizzazione esatta, VQE con termine DM (ansatz $W$-2q.6), dinamica
di Trotter, correlatori dinamici (Hadamard test). \textbf{Non} copre: i
notebook esplorativi del confronto ansatz ($W$ vs RBS-2q --- conclusione
già stabilita, non ripetuta qui: $W$-2q.6 raggiunge $\mathcal F=1$ dove
RBS-2q satura a $\mathcal F\approx0.9999$, vedi
\texttt{log\_decisioni.md}), lo scan sistematico delle 81 combinazioni
correlatore (\texttt{scan81\_*}, catalogazione, non riproduzione
dati/figure), il diagramma schematico del circuito (illustrativo, non
derivato da dati). Questi restano nel Project come documentazione, non
duplicati qui.

\section{Struttura del pacchetto}

\begin{verbatim}
codice/            moduli .py (libreria, mai eseguiti direttamente)
dati/              output .npz/.png degli script di generazione dati
figure/            figure finali (PNG --- vedi nota sotto)
genera_figure/     script di generazione figure
01_..., 05_...     script di dati/validazione, numerati nell'ordine d'uso
esegui_tutto.py    orchestratore
\end{verbatim}

\begin{verbatim}
python3 esegui_tutto.py
\end{verbatim}
Richiede: \texttt{qiskit}, \texttt{numpy}, \texttt{scipy}, \texttt{matplotlib}.

\begin{center}
\renewcommand{\arraystretch}{1.2}
\begin{tabular}{@{}p{4cm}p{6cm}p{4cm}@{}}
\toprule
\textbf{Nota su PNG vs PDF} & \\
\midrule
\multicolumn{3}{p{14cm}}{A differenza del pacchetto di Parte 2, le figure
qui sono \textbf{raster PNG}, non PDF vettoriale --- fedele al
comportamento originale degli script del Project
(\texttt{generate\_figures\_circuito\_trimero\_anello.py} salva solo
\texttt{.png}, mai modificato). Non un difetto introdotto qui.}\\
\bottomrule
\end{tabular}
\end{center}

\section{Le onde di dipendenza reale}

\textbf{Onda 1} (nessun ingresso, indipendenti fra loro): \texttt{01},
\texttt{02}, \texttt{03}, \texttt{04}. \textbf{Onda 2} (dipende da
\texttt{02}): \texttt{05}. \textbf{Onda 3} (dipende da \texttt{04} solo
concettualmente --- la funzione \texttt{classical\_exact} è copiata
localmente, non importata, per evitare un side-effect di \texttt{chdir}):
\texttt{genera\_figure/genera\_figure\_correlatori.py}.

\section{Script 01--05}

\subsection{01\_diagonalizzazione\_esatta.py}
\textbf{Ingresso}: nessuno. \textbf{Uscita}: nessun file dati, ma un
effetto collaterale del self-test: \texttt{dati/trimer\_ring\_exact.png}
(spettro + magnetizzazione, percorso hard-coded nello script originale,
non modificato). \textbf{Chi lo consuma}: nessuno.
\begin{verbatim}
python3 01_diagonalizzazione_esatta.py
\end{verbatim}

\subsection{02\_genera\_vqe\_dm.py}
\textbf{Ingresso}: nessuno (ottimizzazione multistart da zero).
\textbf{Uscita}: \texttt{dati/w2q6\_params\_optimal.npz} (campi:
\texttt{params}, \texttt{E\_vqe}, \texttt{E\_exact}, \texttt{fidelity}).
\textbf{Chi lo consuma}: \texttt{05}, e l'intera pipeline di Parte 2
(pacchetto separato).
\begin{verbatim}
python3 02_genera_vqe_dm.py
\end{verbatim}

\subsection{03\_self\_test\_trotter.py}
\textbf{Ingresso}: nessuno. \textbf{Uscita}: nessun file, solo resoconto
a schermo (4 self-test: coerenza $S_z$, equivalenza circuito/matrice,
convergenza $N\to\infty$, fattorizzazione $H_0$). \textbf{Chi lo
consuma}: nessuno.
\begin{verbatim}
python3 03_self_test_trotter.py
\end{verbatim}

\subsection{04\_valida\_correlatori\_esatti.py}
\textbf{Ingresso}: nessuno (confronto contro benchmark classico
\texttt{scipy.linalg.expm}, calcolato internamente). \textbf{Uscita}:
nessun file, solo resoconto a schermo. \textbf{Chi lo consuma}: nessuno
direttamente --- la funzione \texttt{classical\_exact} è \textbf{copiata}
(non importata) in \texttt{genera\_figure/genera\_figure\_correlatori.py}
per evitare il side-effect di \texttt{chdir} che questo script esegue.
\begin{verbatim}
python3 04_valida_correlatori_esatti.py
\end{verbatim}

\subsection{05\_valida\_correlatori\_vqe.py}
\textbf{Ingresso}: \texttt{dati/w2q6\_params\_optimal.npz} (da
\texttt{02}). \textbf{Uscita}: nessun file, solo resoconto a schermo
(81 combinazioni siti$\times$componenti, preparazione VQE reale vs
ampiezze esatte). \textbf{Chi lo consuma}: nessuno.
\begin{verbatim}
python3 05_valida_correlatori_vqe.py
\end{verbatim}

\section{Script genera\_figure/}

\subsection{genera\_figure\_correlatori.py}
\textbf{Ingresso}: nessuno (calcola tutto da zero: Hamiltoniana esatta +
circuito, stessa logica di \texttt{04} ma con \texttt{classical\_exact}
copiata localmente). \textbf{Uscita}:
\texttt{figure/fig\_trotter\_trimero.png} (convergenza Trotter,
$C_{11}^{yy}(t{=}1.3)$), \texttt{figure/fig\_validazione\_trimero.png}
(validazione continua su un intervallo di $t$),
\texttt{figure/fig\_sito3\_zero\_trimero.png} (zero strutturale del sito
3, $C_{33}^{xz}(t)\equiv0$).
\begin{verbatim}
python3 genera_figure/genera_figure_correlatori.py
\end{verbatim}

\section{Documentazione del codice aggiunta in questo pacchetto}

Le copie di partenza in \texttt{codice/} sono state verificate
bit-per-bit identiche al Project (\texttt{diff}), poi arricchite con
docstring di funzione mancanti (\texttt{vqe\_w2q6\_trimero\_anello.py}:
1/7 $\to$ 7/7; \texttt{trotter\_trimero\_anello.py}: 5/11 $\to$ 11/11),
ciascuna con ruolo nel modulo e ruolo nell'insieme, come richiesto.
Verificato funzionalmente dopo ogni modifica (sintassi + riesecuzione
del self-test corrispondente): nessuna modifica di logica, solo
documentazione.

\section{Sequenza completa (copia-incolla)}

\begin{verbatim}
python3 01_diagonalizzazione_esatta.py
python3 02_genera_vqe_dm.py
python3 03_self_test_trotter.py
python3 04_valida_correlatori_esatti.py
python3 05_valida_correlatori_vqe.py
python3 genera_figure/genera_figure_correlatori.py
\end{verbatim}
Equivalente, in un solo comando: \texttt{python3 esegui\_tutto.py}.

\section{Verifica end-to-end effettuata}

\texttt{dati/} e \texttt{figure/} svuotate, \texttt{esegui\_tutto.py}
rieseguito da zero: \texttt{exit code 0}, nessun \texttt{FALLITO} né
\texttt{Traceback} reale nel log (un falso positivo testuale nel
self-test di \texttt{01}, verificato: fa parte del messaggio esplicativo
del test, non un errore). Numeri rigenerati confrontati con quelli già
noti: $E_\text{VQE}=-5.534370017393404$,
$\mathcal F=0.9999999999996096$ --- coincidenti.

\end{document}
